\label{ch:system-development}

\section{Design Goals and Motivation}

Mlguide was developed as a direct continuation of the work done in the pre-master project \parencite{bergsto2026mlselection}.
The developed roadmap outlines a process for filtering and ranking 
research literature based on user-defined context parameters, including 
ML paradigm, product lifecycle (PLC) phase, and application area. 
This is illustrated with a simplified version of the roadmap in Figure~\ref{fig:ml-roadmap-pre-master}.
In the pre-master project, the main focus was on collecting and analyzing 
research literature. Several directions for future development were identified, 
including better support for problem formulation, ontology-based ML knowledge 
representation, and support for translating recommendations into an 
ML implementation \parencite{bergsto2026mlselection}.

\begin{figure}[H]
    \centering
    \includegraphics[width=0.4\textwidth]{Figures/04-system-design-development/roadmap-pre-master.pdf}
    \caption{Simplified ML selection roadmap from the pre-master project \parencite{bergsto2026mlselection}}
    \label{fig:ml-roadmap-pre-master}
\end{figure}

MLguide addresses these directions by extending the original roadmap into 
a web-based decision-support system. The main design goals were to support 
users in formulating their problem, to recommend relevant ML methods based 
on both structured knowledge and research literature, and to provide a 
starting point for implementation through generated notebooks. 
The system is primarily aimed at engineers and students who are new to 
machine learning and need structured guidance in choosing a suitable method.
A web-based interface was chosen to 
lower the threshold for use, making it easy for other users to access and test
the system without installation. 

\section{Incremental and Iterative Development}

MLguide was developed incrementally over the course of the master's project. 
New features and components were built and deployed one at a time, each adding 
to the working system. This approach is consistent with incremental software 
development, where the system grows through a series of versions rather than 
being built all at once \parencite[33]{sommerville2011software_engineering}.

Figure~\ref{fig:incremental-development} illustrates this process. 
Planning, development, and self evaluation were carried out as concurrent 
activities, producing successive versions of the system. Throughout development, 
the author continuously tested individual features and the overall system output 
to check that behaviour was correct and results were useful.

The development also had an iterative character. Feedback from the thesis 
supervisor informed adjustments to the system design and priorities through the development period. 
In addition, two rounds of user testing were conducted, 
and the findings were used to adjust the system. 
This iterative cycle is consistent with the DSR process described by 
\textcite{peffers2007design_science_research_methodology}, where evaluation 
results inform a return to design and development. The user testing setup is described in the following section.

\begin{figure}[!htbp]
    \centering
    \includegraphics[width=0.8\textwidth]{Figures/04-system-design-development/incremental-development.pdf}
    \caption{Illustration of the incremental development process, based on 
    \textcite[33]{sommerville2011software_engineering}}
    \label{fig:incremental-development}
\end{figure}

\section{User Testing}
\subsection{Setup and Context}

The system was evaluated through two rounds of user testing in a machine 
learning course at the Norwegian University of Science and Technology (NTNU), where students used MLguide as part of their 
course projects. The participants were engineering students where most had limited 
prior experience with machine learning, which made them representative 
of the intended user group. The user testing rounds are summarised in Table~\ref{tab:user-testing}. 

\begin{table}[H]
\centering
\caption{User testing sessions}
\label{tab:user-testing}
\begin{tabular}{l p{5.5cm} p{5.5cm}}
\toprule
 & \textbf{Project 1} & \textbf{Project 2} \\
\midrule
\textbf{Topic} & Anomaly detection in time series data & Reinforcement learning in product lifecycle \\
\addlinespace
\textbf{ML area} & Unsupervised learning & Reinforcement learning \\
\addlinespace
\textbf{Task} & Detect anomalies using unsupervised learning techniques and compare performance & Analyse RL applications across lifecycle phases and implement a sustainability agent \\
\addlinespace
\textbf{MLguide task} & Test the system on their problem and report on relevance and usefulness of results & Assess the applicability of MLguide to their problem \\
\addlinespace
\textbf{Deadline} & March 2026 & April 2026 \\
\bottomrule
\end{tabular}
\end{table}

The two projects differed in scope and focus. In the first project, students 
worked on a well-defined problem: detecting anomalies in time series data using 
unsupervised learning. They were asked to compare at least three techniques, 
which made MLguide a natural fit. Students could use it to get method 
recommendations for their problem, either as inspiration before choosing 
methods or as a comparison after they had already made a selection.

The second project was more open-ended. Students were asked to analyse 
applications of reinforcement learning across four product lifecycle phases: 
product design, manufacturing, product use, and recycling. In addition, 
they were asked to implement a simulated RL environment and a sustainability 
agent for either product design, operations, or product use. 
This made the project less focused on method selection, 
but MLguide could still be used in two ways: to find relevant articles by 
submitting problem descriptions related to different lifecycle phases, 
and to get algorithm recommendations for the implementation part of the project.

\subsection{Feedback and Limitations}

The feedback from both projects gave useful insight into the usability of 
MLguide and what could be improved. However, two limitations apply.

The first concerns the feedback format. Feedback was collected as part of 
the students' project reports, which fit naturally within the course structure 
and allowed students to reflect on MLguide in the context of their own problem 
and data. Since each group worked on a different problem, an open format seemed 
more suitable than a fixed set of questions. The trade-off is that responses 
varied in detail and focus, making it more difficult to draw systematic conclusions.

The second limitation is that the feedback arrived quite late in the development 
period. This was probably unavoidable, as a working system had to be in place 
before it could be tested. The feedback led to some adjustments, but larger 
changes such as new features or structural redesigns were not feasible within 
the remaining time. These can instead be left as directions for future work.

A simple overview of the development timeline, including the user testing sessions, is illustrated in Figure~\ref{fig:development-timeline}.

\begin{figure}[!htbp]
    \centering
    \includegraphics[width=0.8\textwidth]{Figures/04-system-design-development/development_timeline.pdf}
    \caption{Development timeline of MLguide}
    \label{fig:development-timeline}
\end{figure}
